\section{Threats to Validity}
\label{sec:threats}

\paragraph{Internal validity.}
The completeness score formula (Equation~\ref{eq:completeness}) and its component weights
(25/25/20/15/15) were defined a priori based on domain judgment rather than optimized on
a holdout set. Different weight assignments would shift absolute scores but are unlikely to
change the relative ordering of diseases or the qualitative patterns we report.
To verify this, we tested three alternative weight configurations: equal weights
(20/20/20/20/20), a coverage-heavy scheme (35/35/10/10/10), and an accuracy-heavy scheme
(20/20/10/25/25). The disease ranking is preserved across all configurations: Dengue,
Measles, and COVID-19 consistently score highest and Ebola consistently lowest,
confirming that the completeness ordering is robust to reasonable weight variation.
The structural error taxonomy reflects common failure modes observed during development; it may
not cover all possible errors in other disease models.

\paragraph{Construct validity.}
Fuzzy string matching with synonym groups is a heuristic: it may produce false positives
when different epidemiological concepts share a word (e.g., ``Susceptible adults'' matching
``Susceptible children'') or false negatives for non-standard abbreviations not covered by
our synonym list. The 0.6 \texttt{SequenceMatcher} threshold was set to minimize both error
types on our dataset but was not tuned on a separate validation set.

\paragraph{External validity.}
Our study covers 10 diseases with models of varying complexity (4--23 flows). These were
chosen to span a range of standard structures (SIR, SEIR, multi-group, host--vector) but
may not represent the full diversity of published compartmental models, particularly agent-based
or network-based models that do not fit the ODE compartmental structure. Results for
diseases with highly non-standard parameter naming conventions (e.g., Bayesian notation)
may differ from what we report here.
Although evaluated on these 10 ODE compartmental models, the framework is
design-agnostic at the pipeline level: the three-phase structure (reference construction,
LLM-assisted extraction, gap completion) applies to any parametric model class where a
formal reference and a natural-language description coexist, including pharmacokinetic,
ecological, and within-host models.

\paragraph{LLM non-determinism.}
All Phase~2 extractions and Phase~3 inferences used temperature~0, which minimizes but
does not eliminate non-determinism in some LLM backends. A small fraction of results could
differ on replication. To mitigate this, we report averaged patterns and note that all runs
used the same prompts and model versions.

\paragraph{Gold standard quality.}
The reference \texttt{.compmodel} files were constructed from each paper's ODE equations
and parameter tables, and validated against the reported simulation results and figures
in the paper---not authored from informal reading.
Each compartment, flow, and parameter in the gold model traces directly to a named
equation or table entry. For papers that provided simulation code, we cross-validated
the gold model's compartment wiring and initial conditions against the code. For those
without code, ambiguous parameter assignments were resolved by returning to the
differential equation system and identifying which flow each symbol multiplies.
This construction methodology ensures the gold standard is grounded in the paper's own
formal specification, minimising subjective interpretation. The reference models and
their source annotations are included in the project repository.
Any residual misinterpretation of ambiguous notation propagates to the completeness score,
which is itself the instrument for detecting such ambiguity in the source paper.

\section{Conclusion}
\label{sec:conclusion}

We presented a \textbf{measurement framework for epidemiological model completeness} as a
systematic instrument for quantifying the gap between what a paper specifies and what a
reader can reproduce. LLMs were used as simulated reproducers for models specified in epidemiological papers. Phase~1 builds structured reference models
and paper promise inventories. Phase~2 employs LLM-assisted entity extraction from PDFs
across three commercial backends, evaluated with a domain-aware fuzzy string matching
protocol. Phase~3 closes extraction gaps via iterative retrieval-augmented completion,
structural validation, and rule-based repair, then assigns a five-component completeness
score that jointly measures gap reduction, reference agreement, fill traceability,
parameter accuracy, and structural integrity.

Applied to 10 peer-reviewed modeling papers across 10 infectious diseases, the pipeline
achieves a mean completeness score of \textbf{$74.0\%$}. Phase~3 raises mean structural recall
from $0.822$ (Phase~2 draft) to $0.965$ (winner mode), with all three Phase~3 modes improving
recall over the Phase~2 baseline. All 10 diseases show detectable structural errors in
their LLM-extracted models before repair; $44\%$ are resolved automatically.
RAG-only achieves the highest recall in 9 of 10 diseases and is the safest mode because
its fills are traceable to specific paper passages. LLM-only adds genuine value as a
coverage tool, but its fills are unverifiable and lower the traceability component of
the completeness score. The distinction between RAG-sourced and LLM-inferred fills is
itself a measurement outcome: a gap that only LLM can fill---but RAG cannot---is direct
evidence of a specification failure in the paper.

These results quantify a consistent pattern: LLMs reliably extract compartment structure
but struggle with flow paths and parameters embedded in ODE derivations or referenced
only by symbolic expressions. The specification-validation gap is measurable and
non-trivial: even with the enhanced extraction pipeline using GPT-5.4-mini, Phase~2 drafts
miss on average $28\%$ of flows and $12\%$ of parameters relative to the best-case
Phase~2 draft---gaps that Phase~3 reduces substantially, achieving perfect recall for
all three entity types in 7 of 10 diseases.

%\paragraph{Future work.}
%Several concrete directions emerge directly from the failure cases described in this paper.

%\noindent\textbf{Expand the dataset.}
%The most immediate improvement is to include more papers per disease (currently one each).
%A single paper per disease means that a paper-specific writing quirk---parameters in
%figure captions (flu), composite ODE coefficients (TB), Bayesian-only values
%(Ebola)---dominates the per-disease score. Adding two or three papers per disease would
%separate pipeline limitations from individual paper choices, and would enable
%within-disease variance analysis and more robust mean completeness estimates.

%\noindent\textbf{Add figure and table extraction.}
%Several of the worst failures (flu, Ebola) stem from numeric values that exist only
%inside figures, supplementary tables, or figure captions. Integrating OCR-based figure
%parsing and structured table extraction (beyond Camelot's current scope) would close
%this entire class of gaps without any change to the downstream phases.

%\noindent\textbf{Improve RAG for multi-part definitions.}
%TB and Malaria show that parameters defined incrementally across multiple sections are
%not recoverable by single-passage retrieval. A multi-hop RAG strategy---where a
%retrieved passage can trigger further retrieval of referenced symbols---would let the
%system assemble composite definitions that no individual passage contains in full.

%\noindent\textbf{Enhance the evaluation protocol.}
%The current fuzzy string matching works well for standard compartment and flow names
%but still struggles with highly domain-specific notation (e.g., Bayesian posterior
%symbols, host--vector composite terms). Extending the synonym library and adding a
%disambiguation step for homonyms (``Susceptible adults'' vs.\ ``Susceptible children'')
%would improve recall precision and reduce false matches that inflate precision scores.

%\noindent\textbf{Extend the structural validator.}
%Adding dimension-aware validation (e.g., detecting a rate parameter with units
%$\text{person}^{-1}$ assigned to a flow expecting $\text{person}^{-1}\text{day}^{-1}$)
%and a formal parameter-to-flow assignment step (to resolve orphaned parameters in
%multi-group models like HIV and Malaria) would significantly improve the structural
%integrity component, currently the weakest in the completeness score.

%\noindent\textbf{Test generalizability.}
%Applying the framework to stochastic, agent-based, and network models beyond ODE
%compartmental structures would test whether the three-phase approach and the completeness
%score translate to a broader class of mathematical models in infectious disease
%epidemiology. This also includes models where multiple interacting patches or age-strata
%appear as structured population layers rather than separate compartments.

%\noindent\textbf{Close the loop with simulation.}
%Finally, studying whether completeness scores correlate with downstream simulation
%accuracy---i.e., whether a higher-scoring extracted model reproduces the paper's reported
%trajectories more faithfully---would validate the score as a predictor of scientific
%reproducibility and not just a structural inventory metric.
